\scanpage{170}
Comment reconstituer p.ex. les \deleted{ouverts ``globaux''} \deleted{\unsure{par les|pour}} topos canoniques (cf p. 74' \unsure{ff.}) qui sont des ouverts globaux de $X=\mathbb C^2$, \unsure{savoir|à savoir}
\deleted{$\mathbb C^2\setminus\Delta_1$ \unsure{les}}
\[
(278)\quad X=\mathbb C^2,\quad
\underbrace{\mathbb C^2\setminus\Delta_1}_{U_1},\quad
\underbrace{\mathbb C^2\setminus\Delta_2}_{U_2},\quad
U=U_1\cup U_2=\mathbb C^2-\{\mathbb P\},\quad
U_1\cap U_2=\mathbb C^*\times\unsure{\mathbb C^*}
\]
Parmi ceux-ci, \deleted{seuls} le dernier est déjà ``élémentaire'' -- c'est $D_0^*$. On aura
\deleted{$\mathbb C$}
\[
\left.\begin{array}{l}
U_1=\mathbb C\times\mathbb C^*\\
U_2=\mathbb C^*\times\mathbb C
\end{array}\right\}\quad\text{où}\quad
\mathbb C=\mathbb C^*\amalg_{\mathbb D^*}\mathbb D
\]
\[
U=U_1\amalg_{U_1\cap U_2}U_2\qquad\text{où }U_1\cap U_2=D_0^*
\]
\[
X=\mathbb C\times\mathbb C
\]
\deleted{alors il suffit d'exprimer}
$U_1$\inserted{, $U_2$} seront donc donnés par sommes amalgamées
\diagram{\(\displaystyle(279)\quad\left\{\vcenter{\hbox{%
\begin{tikzcd}[ampersand replacement=\&,row sep=large,column sep=large]
\&\mathbb D^*\times\mathbb C^*\arrow[dl]\arrow[dr]\&\\
\mathbb C^*\times\mathbb C^*\arrow[dr,dotted]\&\text{coc.}\&\mathbb D\times\mathbb C^*\arrow[dl,dotted]\\
\& (U_1)\&\\[2em]
\&\mathbb C^*\times\mathbb D^*\arrow[dl,no head]\arrow[dr]\&\\
\mathbb C^*\times\mathbb C^*\arrow[dr,dotted,no head]\&\text{coc.}\&\mathbb C^*\times\mathbb D\arrow[dl,dotted]\\
\&(U_2)\&
\end{tikzcd}}}\right.\)}
\margin{169}
